% Source: Wald GR, physical PDF pp. 16--17
%         (printed pp. 3--4).
% Coverage: Section 1.1, Introduction; no equations.
% Figures/tables/footnotes: none.
% Status: source-reviewed and compile-clean.
% Uncertainties: none.

\subsection{Introduction}\label{sec:1.1}

General relativity is the theory of space, time, and gravitation formulated by
Einstein in 1915. It is often regarded as a very abstruse and difficult theory, partly
because the new viewpoint it introduced on the nature of space and time takes some
effort to get used to since it goes against some deeply ingrained, intuitive notions,
and partly because the mathematics required for a precise formulation of the ideas
and equations of general relativity (namely, differential geometry) is not familiar to
most physicists. Although it has been universally acknowledged as being a beautiful
theory, the potential relevance of general relativity to the rest of physics has not been
universally acknowledged and, indeed, probably for this reason, the subject has lain
nearly dormant during much of its history.

Strong interest in general relativity began to be revived starting in the late 1950s,
particularly by the Princeton group led by John Wheeler and the London group led by
Herman Bondi. Although it is difficult to determine the reasons for trends in
physics, two developments---relating general relativity to other areas of physics and
astronomy---have contributed greatly to the sustained interest in general relativity
which has continued since then. The first is the astronomical discovery of highly
energetic, compact objects---in particular, quasars and compact X-ray sources. It is
likely that gravitational collapse and/or strong gravitational fields play an important
role here, and if so, general relativity would be needed to understand the structure
of these objects. The modern theory of gravitational collapse, singularities, and black
holes was developed beginning in the mid-1960s largely in response to this impetus.

A second factor promoting renewed interest in general relativity is the realization
that although gravitation may be too weak to play an important role in laboratory
experiments in particle physics, nevertheless it is of great importance to our further
understanding of the laws of nature that a quantum theory of gravitation be
developed. In order to make progress toward this goal, a deeper understanding of
some aspects of the classical theory of gravitation---general relativity---may be
needed. Interest in this program has been greatly strengthened by the prediction of
quantum particle creation in the gravitational field of a black hole, as well as by
advances in the study of gauge theories in particle physics.

But even aside from the potential impact of general relativity on astronomy and on
other branches of physics, the theory in its own right makes many remarkable
statements concerning the structure of space and time and the nature of the
gravitational field. After one has learned the theory, one cannot help feeling that one
has gained some deep insights into how nature works.

The purpose of this book is to present the theory of general relativity. We will take
a more modern, geometrical viewpoint than Einstein had, and we will, of course,
discuss the recent advances and developments, but the essential content of the theory
is the one Einstein gave over half a century ago. We begin in this chapter by
discussing the structure of space and time and the basic ideas of relativity theory from
an intuitive, physical point of view. More complete introductory discussions are
given by Geroch (1978a) and Wald (1977a). The remainder of this book will be
devoted to making these ideas mathematically precise and exploring their
consequences.
